\section*{Hoofdstuk \ref{ch:b2:phylogenetic-trees} --- Fylogenetische bomen en de moleculaire klok}

\begin{solution}{exo:b2:phylogenetic-trees:1}
\omterm{def:b2:phylogenetic-trees:tree}{Homologie}: een kenmerk dat van een gemeenschappelijke voorouder is geërfd ---
het opperarmbeen, spaakbeen en de ellepijp van de vleugel van een vleermuis en
de vin van een walvis. \omterm{def:b2:phylogenetic-trees:tree}{Analogie}: een gelijkend kenmerk dat onafhankelijk is
ontstaan --- de vleugels van vleermuizen en vogels als vliegvlakken (de botten
zijn homoloog als voorpoten, de vleugel als zodanig niet), of het cameraoog van
de octopus en van de gewervelden. \omterm{def:b2:phylogenetic-trees:tree}{Synapomorfie}: een afgeleide \omterm{def:b2:phylogenetic-trees:tree}{homologie} die een
\omterm{def:b2:phylogenetic-trees:tree}{clade} deelt en die haar definieert --- veren voor vogels, het amnionei
voor amnioten.
\end{solution}

\begin{solution}{exo:b2:phylogenetic-trees:2}
Vogels: monofyletisch. Reptielen zonder de vogels: parafyletisch.
Vissen zonder de tetrapoden: parafyletisch (longvissen staan dichter bij ons dan bij
forellen). Vliegende gewervelden (vleermuizen, vogels, pterosauriërs):
polyfyletisch. Zoogdieren: monofyletisch. Prokaryoten: parafyletisch (bacteriën
en archaea, met weglating van de eukaryoten --- en archaea staan dichter bij
eukaryoten dan bij bacteriën).
\end{solution}

\begin{solution}{exo:b2:phylogenetic-trees:3}
Zustergroep van de vogels: krokodillen. Zustergroep van de zoogdieren: alle
andere amnioten (schildpadden, hagedissen, krokodillen, vogels samen). Kleinste
\omterm{def:b2:phylogenetic-trees:tree}{clade} met hagedissen en vogels: (hagedissen,(krokodillen, vogels)), de
diapsiden zonder schildpadden op deze boom. Het verwisselen van krokodillen en
vogels op hun knoop verandert niets: een knoop mag vrij worden gedraaid.
\end{solution}

\begin{solution}{exo:b2:phylogenetic-trees:4}
Ongewortelde: $3$; $3\times 5\times 7 = 105$; $3\times 5\times
7\times 9\times 11\times 13\times 15 = \num{2027025}$. Gewortelde bomen
voor 4 taxa: elk van de 3 ongewortelde bomen heeft 5 takken, wat 15 geeft.
\end{solution}

\begin{solution}{exo:b2:phylogenetic-trees:5}
$((A,B),(C,D))$: posities 1, 2, 3 elk één stap, positie 4 twee stappen
($\mathrm{A}$ en $\mathrm{G}$ aan beide kanten), positie 5 één: 6 stappen.
$((A,C),(B,D))$: $2 + 2 + 2 + 1 + 1 = 8$. $((A,D),(B,C))$: $2 + 2 + 2 +
2 + 1 = 9$. De eerste boom is het spaarzaamst; op die boom is positie 4
homoplastisch (de verandering $\mathrm{A}\to\mathrm{G}$, of het omgekeerde,
treedt twee keer op).
\end{solution}

\begin{solution}{exo:b2:phylogenetic-trees:6}
Voeg $A$ en $B$ samen op hoogte 3. Dan $d(AB,C) = 14$, $d(AB,D) = 16$,
$d(C,D) = 10$: voeg $C$ en $D$ samen op hoogte 5. Dan $d(AB,CD) = (14 + 16
+ 14 + 16)/4 = 15$: wortel op hoogte $7.5$. Boom $((A,B),(C,D))$.
\end{solution}

\begin{solution}{exo:b2:phylogenetic-trees:7}
$d = -0.75\ln(1 - 4p/3)$: $0.052$, $0.38$, $1.21$. Bij $p = 0.6$ is het
argument van de logaritme $0.2$ en de helling $\mathrm{d}d/\mathrm{d}p
= 1/(1 - 4p/3) = 5$: een steekproeffout van $0.02$ in $p$ wordt
$0.1$ in $d$, en de aanname van gelijke snelheden op elke positie,
onjuist in elk echt gen, voegt een vertekening van dezelfde grootte toe. Het
gen is verzadigd.
\end{solution}

\begin{solution}{exo:b2:phylogenetic-trees:8}
$t = d/2k = 0.02/(2\times 10^{-9}) = 10$ miljoen jaar. Bij 500
miljoen jaar is $d = 2\times 10^{-9}\times 5\times 10^{8} = 1.0$ en
$p = 0.75(1 - e^{-4/3}) = 0.55$: drie kwart van de weg naar
verzadiging, waar de correctie steil is en de datering onbetrouwbaar
--- voor zulke splitsingen is een trager gen nodig.
\end{solution}

\begin{solution}{exo:b2:phylogenetic-trees:9}
\qty{52}{\%}: de \omterm{def:b2:phylogenetic-trees:tree}{clade} wordt in amper de helft van de hersteekproeven
teruggevonden; de gegevens zeggen er bijna niets over, en de alternatieven
samen zijn even waarschijnlijk. \qty{99}{\%}: het signaal is consistent over
de posities; de \omterm{def:b2:phylogenetic-trees:tree}{clade} is sterk ondersteund (gegeven het model en de
alignering --- systematische fouten zoals \omterm{prop:b2:phylogenetic-trees:pitfalls}{aantrekking van lange takken} worden
door de bootstrap niet opgespoord). Voor de eerste: voeg sequentie toe (meer
genen), voeg taxa toe die de takken onderbreken, en controleer met een andere
methode.
\end{solution}

\begin{solution}{exo:b2:phylogenetic-trees:10}
Twee lijnen die elk op een groot deel van de posities veranderd zijn, hebben
op elke positie waar beide veranderden een kans van één op vier om op dezelfde
base uit te komen (met vier basen en gelijke snelheden). Als elk op
\qty{40}{\%} van de posities is veranderd, zijn ze samen op
\qty{16}{\%} veranderd en komen ze toevallig overeen op \qty{4}{\%} van alle
posities --- wat meer kan zijn dan de fractie posities met de ware
synapomorfieën die elk met zijn echte, traag evoluerende verwanten verbinden.
Spaarzaamheid telt overeenkomsten zonder te vragen waarom ze optreden
en voegt de twee samen. Taxa toevoegen die langs elke lange tak aftakken,
verdeelt die tak in kortere stukken, zodat de toevallige overeenkomsten over
meerdere kortere takken worden verdeeld en de ware synapomorfieën van elk
stuk zichtbaar worden; een likelihoodmodel dat op lange takken veel toevallige
overeenkomsten verwacht, corrigeert er rechtstreeks voor.
\end{solution}

\begin{solution}{exo:b2:phylogenetic-trees:11}
$d = -0.75\ln(1 - 0.12) = 0.096$; $t = d/2k = 0.096/(4\times 10^{-8})
= 2.4$ miljoen jaar. Substituties: $0.096\times \num{16000} =
1500$, relatieve nauwkeurigheid $1/\sqrt{1500} = \qty{2.6}{\%}$: een
statistische fout van $\pm 0.06$ miljoen jaar. Maar de snelheid is
gekalibreerd op andere splitsingen en kan voor deze lijn een factor twee fout
zijn (snelheidsvariatie tussen lijnen, en het verschil tussen de
stamboomsnelheid gemeten over generaties en de snelheid gemeten over miljoenen
jaren); en de divergentie van het gen gaat vooraf aan die van de soorten,
omdat de voorouderlijke populatie polymorf was --- zodat de ware splitsing
recenter is dan die van het gen, met een tijd die van de grootte van de
voorouderlijke populatie afhangt. (De aanvaarde datering, uit veel kerngenen en
fossielen, is 6 tot 7 miljoen jaar: de mitochondriale klok is hier te
snel.)
\end{solution}

\begin{solution}{exo:b2:phylogenetic-trees:12}
\omterm{prop:b2:phylogenetic-trees:pitfalls}{Onvolledige lijnsortering}: wanneer twee soortvormingen dicht op elkaar volgen,
kan een voorouderlijk polymorfisme zo sorteren dat de boom van een gen de
soorten anders groepeert dan de soortenboom; horizontale overdracht: een gen
dat van een andere lijn is ontvangen, draagt de geschiedenis van die lijn;
genduplicatie: een paraloog in de ene soort vergelijken met de ortholoog in
een andere geeft de datering van de duplicatie, niet van de soortvorming. Een
soortenboom verkrijgt men door veel genen uit het hele genoom te gebruiken en
de boom te nemen die de meerderheid ondersteunt (of een model dat een bekende
fractie afwijkende genen verwacht), door orthologen zorgvuldig te kiezen, en
door bij prokaryoten de voorkeur te geven aan ribosomale en informationele
genen die zelden worden overgedragen.
\end{solution}

\begin{solution}{pb:b2:phylogenetic-trees:1}
\textbf{1.} Posities 1, 2 en 4 (toestanden die alleen $W$ en $H$ delen,
afgeleid ten opzichte van de kameel): synapomorfieën van $\{W,H\}$. Posities 3
en 6: synapomorfieën van $\{W,H,C\}$. Positie 5 is een autapomorfie van
$W$, een verandering in één lijn die niets groepeert; geen enkele positie is
homoplastisch op de moleculaire boom.
\textbf{2.} Eén verandering per positie: 6 stappen.
\textbf{3.} Posities 1, 2 en 4 vragen nu elk twee veranderingen, posities 3, 5
en 6 één: 9 stappen.
\textbf{4.} De boom van walvis en nijlpaard, met drie stappen. Drie posities op
duizend is een kleine marge die een paar homoplasieën zouden kunnen
omverwerpen; versterken betekent meer genen, meer taxa (andere
herkauwers, pekari's, herten), en een bootstrap.
\textbf{5.} De \omterm{def:b2:phylogenetic-trees:tree}{clade} komt terug in vijf op zes hersteekproeven: matige
ondersteuning, ruim boven toeval maar onder de \qty{95}{\%}
die gewoonlijk sterk wordt genoemd.
\textbf{6.} De buitengroep moet buiten de groep liggen, maar dichtbij genoeg
dat haar sequentie nog alignerbaar en onverzadigd is; de kameel is een
evenhoevige buiten de groep van rund, varken, nijlpaard en walvis. Een
vis zou op bijna verzadigde niveaus verschillen, zijn toestanden zouden niets
zeggen over welke toestand binnen de zoogdieren voorouderlijk is,
en zijn lange tak zou de snelst evoluerende lijn van de binnengroep
aantrekken.
\textbf{7.} $d = 0.0625$, $0.107$, $0.131$, $0.155$.
\textbf{8.} Voeg $W$ en $H$ samen op hoogte $0.031$. Dan $d(WH,C) =
0.107$, $d(WH,P) = 0.131$, $d(C,P) = 0.131$, alles tot $M$
$0.155$: voeg $WH$ en $C$ samen op $0.054$. Dan $d(WHC,P) = 0.131$: samenvoegen
op $0.065$. Ten slotte $M$ samenvoegen op $0.078$. Boom $(M,(P,(C,(W,H))))$.
\textbf{9.} Ja: dezelfde boom, dezelfde nesting.
\textbf{10.} Een walvis die twee keer zo snel veranderd was, zou verder van
alles afstaan, ook van het nijlpaard; UPGMA, dat elk uiteinde op hoogte nul
plaatst en het dichtste paar samenvoegt, zou eerst nijlpaard met rund
samenvoegen en de walvis erbuiten laten --- precies de boom van de anatomen,
om de verkeerde reden.
\textbf{11.} $0.0625\times 1000 = 62.5$ substituties; standaardafwijking
volgens Poisson $\sqrt{62.5} = 7.9$.
\textbf{12.} $7.9/62.5 = \qty{13}{\%}$.
\textbf{13.} $k = d/2t = 0.107/(1.2\times 10^{8}) = 8.9\times 10^{-10}$
per positie per jaar.
\textbf{14.} $t = 0.0625/(2\times 8.9\times 10^{-10}) = 35$ miljoen
jaar.
\textbf{15.} Varken: $0.131/(1.79\times 10^{-9}) = 73$ miljoen jaar;
kameel: $0.155/(1.79\times 10^{-9}) = 87$ miljoen jaar.
\textbf{16.} $35 \pm 4.5$ miljoen jaar (\qty{13}{\%}).
\textbf{17.} De snelheid schaalt als $1/t_{\text{cal}}$ en de datering als
$t_{\text{cal}}$: $\pm 5/60 = \qty{8}{\%}$, dus $\pm 3$ miljoen jaar
--- samen met de poissonfout ongeveer $\pm 5$ miljoen jaar.
\textbf{18.} $d = -0.75\ln(1 - 0.6) = 0.69$: het gen is in feite op twee
derde van zijn posities veranderd, de correctie heeft $p$ met $1.5$
vermenigvuldigd en haar helling is $2.5$; het gen is voor deze splitsing bijna
verzadigd en zijn datering is weinig waard.
\textbf{19.} De afwezigheid van een kenmerk is geen gedeelde afgeleide toestand:
walvissen verloren de achterpoot en het enkelbot volledig, zodat er niets
te zeggen valt over welk enkelbot ze gehad zouden hebben. Een \omterm{def:b2:phylogenetic-trees:tree}{synapomorfie} is
bewijs voor een \omterm{def:b2:phylogenetic-trees:tree}{clade}; het verlies ervan in één lijn is geen bewijs
tegen het lidmaatschap.
\textbf{20.} De fossiele walvissen hebben het enkelbot met dubbele katrol:
walvissen stammen af van een voorouder die het had, en dat plaatst ze binnen
de evenhoevigen. De moleculaire voorspelling wordt door de gesteenten
bevestigd.
\textbf{21.} Parafyletisch: een voorouder met al zijn afstammelingen
behalve de walvislijn.
\textbf{22.} \omterm{prop:b2:phylogenetic-trees:pitfalls}{Onvolledige lijnsortering} van een voorouderlijk polymorfisme
over twee snelle soortvormingen (rund, nijlpaard en walvis splitsten binnen
enkele miljoenen jaren), of vergelijking van paralogen. Beslis door
veel genen te sequeneren en de boom te nemen die de meeste ervan, en hun
aaneenschakeling, ondersteunen; controleer in elke familie op duplicatie.
\textbf{23.} $d = -0.75\ln(1 - 0.333) = 0.30$; $k = d/2t =
0.30/(7\times 10^{7}) = 4.3\times 10^{-9}$ per positie per jaar, vijf
keer het kerngen.
\textbf{24.} Mitochondriaal DNA wordt gerepliceerd door een polymerase met
een zwakkere proefleescontrole, staat bloot aan de zuurstofradicalen van de
ademhalingsketen, en wordt niet door de nucleaire machinerie hersteld; het
mist ook recombinatie, zodat schade zich opstapelt.
\textbf{25.} Boom $(M,(P,(C,(W,H))))$: walvissen zijn de zustergroep van
de nijlpaarden; $k = 8.9\times 10^{-10}$ per positie per jaar; divergentie van
walvis en nijlpaard $35 \pm 5$ miljoen jaar.
\end{solution}
